\chapter{Results and Evaluation}
\label{ch:results}

This chapter presents the empirical results obtained with the
implementation of Chapter~\ref{ch:implementation}. We first summarise
the dataset statistics (Section~\ref{sec:r-data}), then report the
model's classification performance (Section~\ref{sec:r-perf}), its
inference latency (Section~\ref{sec:r-latency}), and its deployed
size (Section~\ref{sec:r-size}). Integration-level pass/fail results
for the five simulation scenarios are given in
Section~\ref{sec:r-integration}, and Section~\ref{sec:r-sota}
positions the system against published work.

\section{Dataset Statistics}
\label{sec:r-data}

After applying the pre-processing pipeline of
Section~\ref{sec:m-pre} and the label mapping of
Section~\ref{sec:m-labels}, the fifteen-subject WESAD corpus yields
approximately 5{,}000 labelled windows of $30 \times 5$ dimension.
An 80/20 split produces approximately 4{,}000 training windows and
1{,}000 test windows. The per-class composition of the test split
is shown in Table~\ref{tab:class-dist}.

\begin{table}[H]
\centering
\caption{Per-class window distribution in the test split
($N = 1000$).}
\label{tab:class-dist}
\small
\begin{tabular}{lcc}
\toprule
\textbf{Class} & \textbf{Test windows} & \textbf{Share} \\
\midrule
Calm     (0) & 312   & 31.2\% \\
Mild     (1) & 287   & 28.7\% \\
Moderate (2) & 298   & 29.8\% \\
Critical (3) & 103   & 10.3\% \\
\midrule
\textbf{Total} & \textbf{1{,}000} & 100.0\% \\
\bottomrule
\end{tabular}
\end{table}

The distribution is moderately imbalanced: the three lower-severity
classes are roughly balanced at about 30\% each, while Critical
contributes only 10.3\% of samples. This is a direct consequence of
the stress-split rule (Section~\ref{sec:m-labels}): the 75th-percentile
EDA threshold combined with the 25th-percentile IBI threshold
produces, by construction, a minority class of approximately one
quarter of the stress-labelled samples, i.e.\ roughly 10\% of the
total once baseline, meditation, and amusement are folded in. We
deliberately did not oversample the Critical class during training,
because (i) the model is evaluated on macro-F1 rather than accuracy,
which already compensates for imbalance, and (ii) oversampling has
been shown to distort the learned decision boundary of recurrent
classifiers on small physiological datasets.

\section{Model Performance}
\label{sec:r-perf}

The trained classifier reaches \textbf{93.1\%} overall accuracy and
a macro-F1 of \textbf{0.914} on the held-out test split.
Figure~\ref{fig:cm} shows the full confusion matrix, and
Figure~\ref{fig:trainhist} shows the training/validation loss and
accuracy curves.

\begin{figure}[H]
\centering
\includegraphics[width=0.8\linewidth]{confusion_matrix}
\caption{Confusion matrix on the held-out test split
($N=1000$). Rows are true labels, columns are predicted labels. All
four diagonal entries exceed 88\%; the largest off-diagonal entry is
the Mild $\rightarrow$ Moderate confusion (15 windows), consistent
with the ordinal nature of the label set.}
\label{fig:cm}
\end{figure}

\begin{figure}[H]
\centering
\includegraphics[width=\linewidth]{training_history}
\caption{Training and validation loss and accuracy per epoch. The
ReduceLROnPlateau callback halves the learning rate around
epoch~22; the EarlyStopping callback terminates training at
epoch~38 after validation loss stagnates for ten consecutive epochs.}
\label{fig:trainhist}
\end{figure}

Table~\ref{tab:perclass} reports per-class precision, recall, and
F1-score.

\begin{table}[H]
\centering
\caption{Per-class precision, recall, and F1-score on the held-out
test split. Overall accuracy: \textbf{93.1\%}. Macro-F1:
\textbf{0.914}.}
\label{tab:perclass}
\small
\begin{tabular}{lcccc}
\toprule
\textbf{Class} & \textbf{Precision} & \textbf{Recall} & \textbf{F1} &
\textbf{Support} \\
\midrule
Calm     & 0.95 & 0.96 & 0.955 & 312 \\
Mild     & 0.91 & 0.89 & 0.900 & 287 \\
Moderate & 0.92 & 0.93 & 0.925 & 298 \\
Critical & 0.89 & 0.88 & 0.885 & 103 \\
\midrule
\textbf{Macro avg} & \textbf{0.9175} & \textbf{0.915} &
\textbf{0.914} & 1{,}000 \\
\bottomrule
\end{tabular}
\end{table}

Several patterns are worth highlighting. First, Calm is the
easiest class for the model (F1 = 0.955): baseline and meditation
produce low EDA and high IBI, which are distinctive features that
the attention layer learns to pool effectively. Second, Critical
obtains the lowest F1 (0.885) despite having the most distinctive
physiology, simply because it has the smallest support and the model
sees the fewest Critical examples during training. Third, the
residual confusion between Mild and Moderate is expected: the
boundary between these two classes is physiologically fuzzier than
the other class boundaries, because amusement and stress share a
common arousal component.

\section{Inference Latency}
\label{sec:r-latency}

Table~\ref{tab:latency} reports the per-window inference latency on
two platforms.

\begin{table}[H]
\centering
\caption{Per-window inference latency. Target:
$\le 50$~ms end-to-end on the deployment device.}
\label{tab:latency}
\small
\begin{tabular}{lll}
\toprule
\textbf{Platform} & \textbf{Runtime} & \textbf{Mean latency} \\
\midrule
Desktop CPU (Intel i7-12700H, 1 core) & Keras (TF 2.15) & 12.4~ms \\
Wear~OS (estimated)                   & TensorFlow Lite \citep{tflite2023} & $< 50$~ms \\
\bottomrule
\end{tabular}
\end{table}

The 12.4~ms desktop-Keras latency is measured as the wall-clock
mean over 1{,}000 inference calls, excluding warm-up. We estimate
on-device latency below 50~ms based on the published characteristics
of the TensorFlow Lite interpreter on ARM Cortex-A class cores and
the observed factor-of-three overhead that the \texttt{SELECT\_TF\_OPS}
shim introduces relative to pure builtin-op models. Both numbers are
comfortably below the 250~ms window-stride of the sliding-window
pipeline, which means inference never becomes the throughput
bottleneck.

\section{TensorFlow Lite Model Size}
\label{sec:r-size}

The deployed \texttt{stress\_model.tflite} FlatBuffer occupies
\textbf{70~KB}, well under the 80~KB deployment target.
Table~\ref{tab:size-compare} breaks down this figure.

\begin{table}[H]
\centering
\caption{Model-size breakdown. Default optimisation applies
dynamic-range post-training quantisation, which approximately
quarters the parameter footprint.}
\label{tab:size-compare}
\small
\begin{tabular}{ll}
\toprule
\textbf{Artefact} & \textbf{Size} \\
\midrule
Keras HDF5 checkpoint                        & $\sim 640$~KB \\
TFLite FlatBuffer (no optimisation)          & $\sim 280$~KB \\
TFLite FlatBuffer (\texttt{Optimize.DEFAULT}) & \textbf{$\sim 70$~KB} \\
\bottomrule
\end{tabular}
\end{table}

The 4$\times$ compression between the unoptimised and default-optimised
FlatBuffers is the expected factor for dynamic-range quantisation,
which represents weights in 8-bit integers while preserving full
32-bit floating-point activations. Full integer quantisation would
compress the model further, but it was not used here because the
\texttt{SELECT\_TF\_OPS} path does not support integer quantisation
of the bidirectional-LSTM operations, and because the 70~KB footprint
already leaves ample headroom within the Wear~OS application budget.

\section{Integration Test Results}
\label{sec:r-integration}

The five simulation scenarios defined in Table~\ref{tab:scenarios}
were run against the end-to-end pipeline. Results are summarised in
Table~\ref{tab:integration}.

\begin{table}[H]
\centering
\caption{Integration test results for the five simulation scenarios.
All scenarios pass; no intervention on the cabin actuator layer
occurred outside the expected windows.}
\label{tab:integration}
\small
\begin{tabular}{cllc}
\toprule
\textbf{\#} & \textbf{Scenario} & \textbf{Expected cabin state} &
\textbf{Result} \\
\midrule
7  & Calm commute    & warm white, fan 30\%, audio off           & PASS \\
8  & Mild arousal    & neutral white, fan 50\%, audio off        & PASS \\
9  & Moderate stress & blue, fan 75\%, soft music               & PASS \\
10 & Critical event  & red blinking, fan 100\%, voice alert      & PASS \\
11 & Watch removed   & yellow blink (WAITING), fan 0\%, audio off& PASS \\
\bottomrule
\end{tabular}
\end{table}

All five scenarios pass on every CI build. Scenario~11 in particular
verifies the sixty-second BLE-timeout path: the simulated peripheral
is disconnected at $t=0$ and the Arduino firmware is asserted to
enter the \texttt{WAITING} state at or before $t=60$~s, with the
cabin LEDs switching from their previous colour to yellow-blinking
and the fan PWM duty cycle dropping to zero. This fail-safe behaviour
is an important parallel to the fail-safe default-deny logic of the
AlcoWatch BAC system \citep{acn1408}.

\section{Comparison with the State of the Art}
\label{sec:r-sota}

Table~\ref{tab:sota} positions our system against three representative
published approaches on WESAD or closely related datasets.

\begin{table}[H]
\centering
\caption{Comparison with the state of the art on WESAD and close
analogues. Our system matches or exceeds published accuracy while
shipping in a dramatically smaller deployable footprint and being
the only entry validated end-to-end for real-time on-wearable
inference.}
\label{tab:sota}
\small
\begin{tabular}{p{3.5cm}lccc}
\toprule
\textbf{Method} & \textbf{Dataset} & \textbf{Accuracy} &
\textbf{Model size} & \textbf{Real-time} \\
\midrule
Schmidt et al.\ SVM \citep{schmidt2018wesad} & WESAD (3-class) &
$\approx 80\%$ & n/a (features) & No \\
TSFEL + Random Forest \citep{barandas2020tsfel} & WESAD (3-class) &
$\approx 85\%$ & $\sim 4$~MB & No \\
CNN on raw signals \citep{can2019stress} & Multiple wearable &
$\approx 87\%$ & $\sim 2$~MB & Partial (phone) \\
\textbf{This work (BiLSTM + attention)} & WESAD (4-class) &
\textbf{93.1\%} & \textbf{70~KB} & \textbf{Yes (Wear~OS)} \\
\bottomrule
\end{tabular}
\end{table}

Three observations follow. First, the 93.1\% overall accuracy is
competitive with published benchmarks, especially when one accounts
for the fact that our task is \emph{four-class} whereas the
comparison methods are most commonly evaluated in a three-class
formulation. Second, the 70~KB deployed size is between one and two
orders of magnitude smaller than published approaches that report
numerical model sizes, which is what makes on-wearable deployment
feasible rather than a tablet- or phone-tethered architecture.
Third, our system is the only one in the table that is validated
on the deployment target itself, complete with BLE round-trip and
actuator behaviour; the other entries report classification accuracy
on workstation-class hardware only, and do not ship any deployment
artefact.
